%% kapa: exam=2 round=19 year=2008 period=3 subject=compensation_law subject_ko=감정평가 및 보상법규 seq=2 number=2 points=40
%% source: Q-Net 감정평가사 기출문제 (cst00309) artlSeq=5024177 file=extracted/1008/5024177/2008년 제19회 감정평가사 2차시험 기출문제/3교시.hwp
%% conversion: kordoc 4.24.51 -> tools/gen2.py (auto); review=proofread-sonnet; flags=none
\begin{longquestion}{【문제 2】}
토지에 대한 개별공시지가결정을 다투려고 하는 경우 다음 각각의 사안에 대하여 논술하시오. (40점) \par
\medskip\noindent (1) 甲은 A시장이 자신의 소유 토지에 대한 개별공시지가를 결정함에 있어서 부동산 가격공시 및 감정평가에 관한 법률 제9조 제2항에 의하여 국토해양부장관이 작성한 토지가격비준표를 고려하지 않았다고 주장한다. 이에 A시장은 토지가격비준표를 고려하지 않은 것은 사실이나, 같은 법 제11조 제4항의 규정에 따른 산정지가검증이 적정하게 행해졌으므로, 甲 소유의 토지에 대한 개별공시지가결정은 적법하다고 주장한다. A시장 주장의 타당성에 대하여 검토하시오. (20점) \par
\medskip\noindent (2) 乙은 A시장이 자신의 소유 토지에 대한 개별공시지가를 결정함에 있어서 부동산 가격공시 및 감정평가에 관한 법률 제11조 제4항에 의하여 받아야 하는 산정지가검증을 거치지 않았다는 이유로 개별공시지가결정이 위법하다고 주장하였다. A시장은 乙의 주장이 있자 산정지가검증을 보완하였다. 乙이 검증절차의 위법을 이유로 개별공시지가결정을 다투는 소송을 제기하려는 경우 그 방법 및 인용가능성은? (20점) \par
\end{longquestion}
